\section{Jensen 不等式与随机凸性}
\label{sec:05-Probability/07-Transforms-and-Bounds/03}

一个日志清洗任务每天要处理固定的 \(1\) TB 数据。集群状态好的时候吞吐 \(300\) MB/s，状态差的时候只有 \(100\) MB/s，两种情况各占一半。平均要跑多久？

顺手的算法是先平均吞吐：\((300+100)/2=200\) MB/s，\(10^6\) MB 除以 \(200\)，得 \(5000\) 秒。

但真实的两天是这样过的：好的那天 \(10^6/300\approx3333\) 秒，差的那天 \(10^6/100=10^4\) 秒，平均 \(6667\) 秒。比刚才那个数多出三分之一，而且方向固定——不管吞吐怎么波动，先平均再折算总是偏小。

出错的那一步是把期望搬过了 \(1/x\)：

\[
\frac{1}{\E[X]}\ne\E\Bigl[\frac{1}{X}\Bigr].
\]

两边不等并不稀奇，稀奇的是不等号方向从不摇摆。为什么低吞吐那天对平均耗时的拉扯，总是压过高吞吐那天的补偿？这需要一个关于函数弯曲方向的结论，而不是对某个具体分布的计算。

\subsection{弦在图像上方，平均就穿不过去}

先把问题剥到最简单的形状。设 \(X\) 以各半概率取 \(x_1\) 和 \(x_2\)，考察一个函数 \(g\)。\(g(\E X)\) 是曲线在两点中点处的高度，而 \(\E[g(X)]\) 是两个端点函数值的平均，也就是连接两点的那条弦在中点处的高度。两个数落在同一条竖线上，谁高谁低完全由曲线弯向哪边决定。

\begin{center}
\begin{tikzpicture}[x=1.5cm,y=0.45cm,font=\footnotesize,>=stealth]
  \draw[->,black!70] (0.2,0) -- (3.95,0);
  \draw[->,black!70] (0.3,0) -- (0.3,13.4);
  \draw[blue!65,thick] plot[domain=0.35:3.55,samples=60] (\x,{\x*\x});
  \node[right,blue!65] at (3.5,12.6) {\(g\) 凸};
  \draw[black!70,thick] (0.6,0.36) -- (3.4,11.56);
  \node[above left,black!70] at (2.9,8.4) {弦};
  \fill[black!70] (0.6,0.36) circle (1.5pt);
  \fill[black!70] (3.4,11.56) circle (1.5pt);
  \draw[black!40,dashed] (2,0) -- (2,3.9);
  \fill[orange!85!black] (2,5.96) circle (1.8pt);
  \fill[blue!70] (2,4) circle (1.8pt);
  \draw[<->,black!60] (2,4.2) -- (2,5.76);
  \node[right,black!70] at (2.08,5.2) {Jensen 间隙};
  \node[left,orange!85!black] at (1.9,6.5) {\(\E[g(X)]\)};
  \node[below right,blue!70] at (2.05,3.9) {\(g(\E X)\)};
  \node[below] at (0.6,-0.1) {\(x_1\)};
  \node[below] at (2,-0.1) {\(\E X\)};
  \node[below] at (3.4,-0.1) {\(x_2\)};
\end{tikzpicture}

\emph{\(X\) 各半概率取 \(x_1,x_2\)。先映射再平均落在弦上，先平均再映射落在曲线上。凸函数的弦压在图像上方，于是 \(\E[g(X)]\ge g(\E X)\)，两者之差就是波动被凸性收取的``手续费''。}
\end{center}

这张图就是 Jensen 不等式的全部几何。把两点换成任意分布，把等权换成任意概率权重，结论不变：

\begin{theorem}[Jensen 不等式]
若 \(g\) 在 \(X\) 的取值范围上是凸函数，且相关期望存在，则 \(g(\E X)\le\E[g(X)]\)。\(g\) 为凹函数时不等号反向。
\end{theorem}

凸性的定义正是``弦不低于图像''这句话的代数版本，\(g(\lambda x+(1-\lambda)y)\le\lambda g(x)+(1-\lambda)g(y)\)。关于凸函数的完整刻画与判别方法，见\hyperref[sec:09-Convex-Optimization/03-Convex-Functions/01]{``弦在图像上方意味着什么''}；凹函数只是凸函数取负，所有结论对称照搬，不必另记一套。

线性函数是唯一穿得过去的一类：\(\E[aX+b]=a\E[X]+b\)。期望本身是线性算子，它与线性运算交换毫无阻力；一旦函数弯了，先后顺序就开始要价，而凸性决定了价钱收在哪一边。

\subsection{一条切线就能证完}

若 \(g\) 可微，凸性还有一个等价说法：图像永远在任一点的切线之上，\(g(x)\ge g(m)+g'(m)(x-m)\)。取 \(m=\E X\)，把 \(x\) 换成随机的 \(X\)，两边取期望，右端的线性项因 \(\E[X-\E X]=0\) 整个消失，剩下的正是 Jensen。

三行的分量不轻。切线是 \(g\) 在均值处的线性化，凸性保证线性化只会低估真值；而线性化的平均恰好等于均值处的函数值，因为线性函数与期望可以交换。两边一对比，不等号就落定了。\(g\) 不可微时（例如 ReLU 在原点的折角）把导数换成次梯度，同一段话原样成立。

取等的条件比``\(X\) 是常数''稍宽一点。若 \(g\) 严格凸，图像上没有任何直线段，那么分布只要散在两个不同点上，弦就严格高于曲线，间隙为正。但若 \(g\) 在某段区间上恰好是仿射的，而 \(X\) 的支撑整个落在那段里，即使 \(X\) 不是常数也会取等。准确的说法是：等号成立当且仅当 \(g\) 在 \(X\) 支撑的凸包上几乎必然是仿射的。

\subsection{对号入座：三个反复出现的方向}

平方函数 \(g(x)=x^2\) 在整条实轴上严格凸，Jensen 给出 \(\E[X^2]\ge(\E X)^2\)，移项就是 \(\Var(X)\ge0\)。方差非负这件事平常没人往 Jensen 上想，但它确实是同一张图：方差正是平方函数收取的那笔手续费。

倒数函数 \(g(x)=1/x\) 在 \(x>0\) 上严格凸，于是 \(\E[1/X]\ge1/\E[X]\)，开头那个多出三分之一的耗时到此有了名字。它的适用面比集群任务宽：固定大小文件的传输时间对带宽波动、固定工作量的完成时间对效率波动、并联电阻的等效阻值对阻值波动，都是同一个不等号。共同点是分子固定而分母随机——低值端对结果的杠杆远大于高值端，波动因此系统性地推高平均。要留神的是``分子固定''这个前提：如果每天的数据量本身也随吞吐一起变，比较的就不再是 \(1/\E X\) 与 \(\E[1/X]\)，Jensen 管不到。

对数函数 \(g(x)=\log x\) 在 \(x>0\) 上严格凹，方向翻过来：

\[
\E[\log X]\le\log\E[X].
\]

平均的对数不小于对数的平均。这个不等号是变分推断和 EM 算法的整个来源。含隐变量 \(Z\) 的模型里，对数似然写成 \(\log\E_{q}[\,p(x,Z)/q(Z)\,]\) 的形式后直接卡住——对数套在期望外面，没法拆开求导。用一次 Jensen 把对数换进期望里面，得到的下界就是证据下界（ELBO），此后优化的对象从不可积的似然变成了可采样的期望，而 Jensen 间隙恰好等于近似后验与真后验之间的 KL 散度。EM 算法的每一步做的是同一件事：造一个贴住当前点的下界，把它最大化，再重新造。细节见\hyperref[sec:13-Machine-Learning/06-Probabilistic-Models/04]{``EM 用下界交替完成隐变量学习''}与\hyperref[sec:14-Deep-Learning/06-Variational-Autoencoders/01]{``VAE 把不可积后验变成可训练下界''}。

同一个不等号还顺手给出了一条经典结果。令 \(X\) 等概率取正数 \(x_1,\ldots,x_n\)，对 \(\log\) 写 Jensen 再两边取指数：

\[
\Bigl(\prod_{i=1}^nx_i\Bigr)^{1/n}\le\frac1n\sum_{i=1}^nx_i.
\]

几何平均不超过算术平均。取 \(2\) 和 \(8\) 验一下，几何平均 \(4\)，算术平均 \(5\)。被独立发现过许多次的算术--几何均值不等式，在这里只是 Jensen 配上一个对数。

日常汇报里同一个陷阱换着面孔出现：多个分片上的困惑度（perplexity）不能先平均再取指数、多次实验的加速比不能直接算术平均、几段路程的平均速度不是各段速度的平均。判断方法始终一样——先问要报告的量是 \(g(\E X)\) 还是 \(\E[g(X)]\)，再看 \(g\) 弯向哪边。

\subsection{方向由函数决定，不是由直觉决定}

Jensen 用错的方式基本只有三种，都出在检查环节。

第一种是函数在支撑上没有一致的弯曲方向。\(x^3\) 在正半轴凸、负半轴凹，拐点在零。若 \(X\) 跨过零点，\(\E[X^3]\) 与 \((\E X)^3\) 的大小关系取决于分布往哪边偏，Jensen 不给结论。这类函数只能把定义域切成凸段与凹段分别处理。

第二种是期望不存在。若 \(\E[g(X)]=\infty\)，不等式形式上仍然成立，却不再携带任何定量信息；重尾分布配上凸函数很容易落进这一格。

第三种是取值跑出定义域。\(\log X\) 要求 \(X>0\) 几乎必然，\(1/X\) 要求 \(X\) 几乎必然不为零——只要有一小块概率质量越界，期望可能根本没有定义。实践中这条最容易被忽略：一个理论上恒正的量，在数值实现里因为下溢或截断出现了零，整个论证随之作废。

条件版本几乎不需要额外说明：\(g(\E[X\given Y])\le\E[g(X)\given Y]\) 几乎必然成立，两边再取一次全期望就回到无条件版本。它是条件方差分解、Rao--Blackwell 定理和条件熵性质里共同的那一步，出场频率比无条件版本还高。

上一节那张图里最下面一条曲线要用到 \(e^{tX}\)，而指数函数是严格凸的，Jensen 立刻给出 \(\E[e^{tX}]\ge e^{t\E[X]}\)：矩母函数不会低于均值处的指数值。下一节把这件事倒过来用——先用指数函数把偏差放大，再用 Markov 把放大后的量压住，最后挑一个最合适的放大倍率。凸性在那里不再是被观察的性质，而是被驱使的工具。
